\section{模型的评价与推广}

\subsection{模型优点}

\begin{enumerate}
  \item \textbf{框架统一、可递进。} 以“情景生成 + 两阶段随机规划 + 执行层闭环”为主线，四个问题只改变不确定性来源，物理约束始终一致，便于横向比较。
  \item \textbf{决策口径严谨。} 情景只取决策日之前的历史日，第一阶段变量全情景共享，共同出发点约束封堵起点套利，扰动检验表明无前视泄漏。
  \item \textbf{对冲与校正并重。} 报童下界抑制补救层的乐观偏差，终端价值 $\theta$ 维持隔夜缓冲，充电来源约束消除紧急电套利。
  \item \textbf{结果可信、对照量化。} 各项残差仅 $10^{-9}$ 量级，无越限与同时充放电；相对无储能方案，购电费用降低 $26.92\%$。
  \item \textbf{含反直觉发现。} 更新预报并非越频繁越好（6:00 边际收益为负），价格信息的价值随计划保守程度递减，均有对照实验支撑。
\end{enumerate}

\subsection{模型缺点}

\begin{enumerate}
  \item \textbf{情景池受历史覆盖限制。} 候选池的月份范围决定精度，负载季节跃升时会触发紧急购电；核加权使有效情景数偏少，可能低估尾部风险。
  \item \textbf{保守设定代价明显。} 第二问模型全年弃置电量占计划购电量的 $17.61\%$，报童下界取 $0.8$ 分位使总费用上升约 $19.5\%$，偏保守而非最优。
\end{enumerate}

\subsection{模型的改进方向}

针对上述缺点，可从以下两方面改进：

\begin{enumerate}
  \item \textbf{提高情景池的覆盖能力。} 引入季节跃升检测，在负载或净负荷突变时动态扩大候选池的月份环形距离或叠加趋势项，避免季节跃升边界出现紧急购电；用 CVaR 或更重尾的情景构造补充核加权对尾部风险的低估。
  \item \textbf{以费用口径自适应选定保守参数。} 报童下界分位数 $q$ 与核带宽 $h$ 目前分别独立选取，可在统一的费用口径下交叉验证、联合搜索，用实际结算费用作为选参标准，避免固定取 $0.8$ 分位带来的过度保守。
\end{enumerate}

\subsection{模型的推广}

本文提出的“情景生成 + 两阶段随机规划 + 情景 MPC + 终端价值 + 报童下界”框架，本质上是含储能设备、面临多重时序不确定性、需要日内滚动决策的联合调度问题的一般解法。这一思路可直接迁移到含储能的虚拟电厂调度、风光储综合能源系统运行、用户参与电力现货市场与需求响应等场景：只需把“负载—光伏”的不确定性替换为风电或价格等随机量，把“购电 + 调整 + 紧急购电”的结算规则替换为对应市场的出清与偏差考核规则，即可复用本模型。这使本文工作不仅回答题目四问，也为电力市场环境下的储能经济调度提供了可落地的建模范式。
